% Generated draft: requires mathematical and source review.
\documentclass[11pt]{article}
\usepackage[margin=25mm]{geometry}
\usepackage{fontspec}
\setmainfont{Latin Modern Roman}
\usepackage{amsmath,amssymb,mathtools,unicode-math}
\setmathfont{Latin Modern Math}
\usepackage{enumitem,needspace,xurl,hyperref}
\setlength{\parindent}{0pt}
\setlength{\parskip}{5pt}
\setlength{\emergencystretch}{3em}
\newcommand{\N}{\mathbb{N}}
\newcommand{\Z}{\mathbb{Z}}
\newcommand{\Q}{\mathbb{Q}}
\newcommand{\R}{\mathbb{R}}
\newcommand{\C}{\mathbb{C}}
\newcommand{\F}{\mathbb{F}}
\newcommand{\A}{\mathbb{A}}
\newcommand{\PP}{\mathbb{P}}
\newcommand{\E}{\mathbb{E}}
\newcommand{\eps}{\varepsilon}
\newcommand{\dd}{\,\mathrm{d}}
\begin{document}
\section*{UnsolvedMath ID: 30004652}
\textbf{Source problem code:} catalogue entry
\subsection{Definitions and mathematical statement}
For terminology, a hyperbolic curve over a field \(k\) conventionally means a smooth geometrically connected curve \(C\) whose smooth projective compactification has genus \(g\) and whose geometric complement consists of \(r\) points, with \(2g-2+r>0\). Thus, in genus zero, \(r\geq 3\). A hyperbolic polycurve conventionally admits a tower
\[
X=X_m\longrightarrow X_{m-1}\longrightarrow\cdots\longrightarrow X_0=\operatorname{Spec}k
\]
of smooth morphisms whose geometric fibers are hyperbolic curves. These conventions explain the terminology; the supplied record does not specify the field or the precise permitted class of towers.

For an integer \(n\geq 1\), the ordered configuration space of \(C\) is
\[
\operatorname{Conf}_n(C)=C^n\setminus\bigcup_{1\leq i<j\leq n}\Delta_{ij},
\]
where products are over \(k\) and \(\Delta_{ij}\) is the locus where the \(i\)-th and \(j\)-th coordinates coincide. The unordered configuration space is the quotient by the symmetric group \(S_n\). The source does not specify which convention is intended or the allowed values of \(n\).

For a connected scheme \(Z\) with geometric base point \(\bar z\), write \(\pi_1^{\mathrm{\acute et}}(Z,\bar z)\) for its profinite \'{e}tale fundamental group. The supplied reconstruction question has the schematic form
\[
\pi_1^{\mathrm{\acute et}}(X,\bar x)\cong
\pi_1^{\mathrm{\acute et}}(Y,\bar y)
\quad\Longrightarrow\quad X\cong Y,
\]
where \(X\) is a hyperbolic polycurve and \(Y\) is a configuration space of a hyperbolic curve, subject to the source's qualification ``in genus zero.''

\textit{Unresolved scope.} This is a schematic question, not a fully quantified conjecture. The record does not determine the base field, whether genus zero refers to the curve defining \(Y\), the fibers in a polycurve tower, or both, or whether the groups are arithmetic or geometric. It also does not specify whether the group isomorphism must be continuous and compatible with additional arithmetic structure, or whether the resulting variety isomorphism is required over \(k\) or over an algebraic closure. These choices must be recovered from the source before a precise universal statement can be asserted. The stated conclusion asks for existence of a variety isomorphism; the record does not explicitly require it to induce the given group isomorphism.
\subsection{Short English statement}
The supplied question asks whether \ a hyperbolic polycurve and a configuration space of a hyperbolic curve must be isomorphic when their \'{e}tale fundamental groups are isomorphic, under an unspecified genus-zero restriction. The field, configuration convention, arithmetic compatibility, and meaning of genus zero require clarification. The supplied status is \texttt{partially\_solved}; this draft does not independently verify that assessment.
\subsection{Sources}
\begin{itemize}
\item Supplied local UnsolvedMath registry record, permanent ID \(30004652\), source code \texttt{catalogue entry}; local registry paths \texttt{data/unsolvedmath/problems.json} and \texttt{catalog/problems.json}. Title: \emph{Anabelian Reconstruction of Genus-Zero Configuration Spaces}. The record attributes the question to the Oberwolfach Reports collection and supplies proposed year 2021, but names no individual proposer.
\item Supplied report citation: \emph{Homotopic and Geometric Galois Theory} (2022), Oberwolfach Reports, pp.~663--744, DOI \texttt{10.4171/owr/2021/12}. \url{https://doi.org/10.4171/owr/2021/12}. The literature list also labels this reference ``Oberwolfach Report 12/2021.''
\item Supplied literature reference: ``Mochizuki, refined genus-zero Grothendieck conjecture.'' \url{https://arxiv.org/abs/2407.09906}.
\item Supplied literature reference: ``Mochizuki, hyperbolic polycurves.'' \url{https://arxiv.org/abs/1902.02058}.
\item Literature assessment dated exactly \texttt{2026-08-22}: the supplied curation labels the problem partially solved and states that Grothendieck-type reconstruction is proved for several classes of hyperbolic polycurves and configuration spaces, while the genus-zero configuration-space case in the stated full generality remains unresolved. This is supplied dated curation derived from the OWR material, not an independently verified literature assessment; the record explicitly requests independent verification of time-sensitive status and cited claims.
\end{itemize}
\end{document}
